\documentclass[11pt,a4paper]{article}

% ---------- Packages ----------
\usepackage[utf8]{inputenc}
\usepackage[T1]{fontenc}
\usepackage{lmodern}
\usepackage[margin=1in]{geometry}
\usepackage{graphicx}
\usepackage{booktabs}
\usepackage{tabularx}
\usepackage{array}
\usepackage{longtable}
\usepackage{xcolor}
\usepackage{listings}
\usepackage{hyperref}
\usepackage{cite}
\usepackage{float}
\usepackage{enumitem}
\usepackage{caption}

% ---------- Hyperref setup ----------
\hypersetup{
    colorlinks=true,
    linkcolor=blue!60!black,
    citecolor=blue!60!black,
    urlcolor=blue!60!black,
    pdfauthor={Mohamed Marzouk et al.},
    pdftitle={VISUALine: A Modular System for Local AI Video Processing on Consumer GPUs}
}

% ---------- Table helpers ----------
\newcolumntype{Y}{>{\raggedright\arraybackslash}X}
\newcolumntype{L}[1]{>{\raggedright\arraybackslash}p{#1}}
\newcolumntype{R}[1]{>{\raggedleft\arraybackslash}p{#1}}

% ---------- Code listing style ----------
\definecolor{codebg}{RGB}{248,248,248}
\definecolor{codeframe}{RGB}{220,220,220}
\definecolor{codekeyword}{RGB}{0,80,160}
\definecolor{codestring}{RGB}{120,60,0}

\lstdefinestyle{yamlstyle}{
    backgroundcolor=\color{codebg},
    frame=single,
    rulecolor=\color{codeframe},
    basicstyle=\ttfamily\small,
    keywordstyle=\color{codekeyword}\bfseries,
    stringstyle=\color{codestring},
    breaklines=true,
    columns=fullflexible,
    keepspaces=true,
    showstringspaces=false,
    tabsize=2
}

\lstdefinestyle{pythonstyle}{
    backgroundcolor=\color{codebg},
    frame=single,
    rulecolor=\color{codeframe},
    basicstyle=\ttfamily\small,
    keywordstyle=\color{codekeyword}\bfseries,
    stringstyle=\color{codestring},
    breaklines=true,
    columns=fullflexible,
    keepspaces=true,
    showstringspaces=false,
    tabsize=4,
    language=Python
}

% ---------- Title ----------
\title{\textbf{VISUALine: A Modular System for Local AI Video Processing on Consumer GPUs}}
\author{
Mohamed Marzouk \and Theodoros Milad \and Yahia \and Mohamed Ezzat \and Kareem Moamen \and Mazen Gaber\\
\small Faculty of Computers and Artificial Intelligence, Capital (formerly Helwan) University, Egypt
}
\date{}

\begin{document}
\maketitle

\begin{abstract}
Local AI video processing is a challenging and highly practical problem. Although recent years have witnessed remarkable progress in computer vision models for detection, segmentation, super-resolution, frame interpolation, and enhancement, there is relatively less work on combining these models into complete local workflows. Existing tools often depend on cloud services, fixed pipelines, or isolated model demos, and they do not fully address issues such as GPU memory usage, video I/O, audio preservation, and reusable user controls. We address this problem from a systems perspective by proposing \textbf{VISUALine}, a modular local AI video and image processing system for consumer GPUs. VISUALine represents workflows as configurable pipelines composed of reusable processing nodes. These pipelines support tasks such as prompt-based cinematic blur, privacy redaction, old-video enhancement, smart vertical reframing, and slow-motion interpolation. Furthermore, VISUALine introduces a shared RGB/RGB+Mask convention, GPU-aware resource management, FFmpeg-based video handling, and a YAML-driven workflow interface. This paper presents VISUALine as an applied system for composing, running, and evaluating practical local AI video-processing workflows.
\end{abstract}

\noindent\textbf{Keywords:} local AI, video processing, computer vision systems, consumer GPU, modular pipelines, prompt-based editing, video enhancement, desktop AI application

% ============================================================
\section{Introduction}

Recent computer vision models can already perform many operations that are useful in video editing and visual media processing. A detector can locate a subject from a text prompt, a segmentation model can separate the subject from the background, a super-resolution model can improve low-resolution frames, and an interpolation model can generate additional frames for slow-motion output~\cite{liu2023groundingdino,ravi2024sam2,wang2021realesrgan,huang2022rife}. These models are usually presented as separate research methods or standalone demos. However, a user-facing video-processing application requires more than a successful model call on a single frame.

A complete local workflow has to solve several practical problems at the same time. It must decode video files, preserve the original audio, move data between different model formats, process long frame sequences without exhausting GPU memory, and write a playable output file. It also has to expose useful controls without forcing the user to understand every internal model parameter. These details are not the main focus of most model papers, but they strongly affect whether the model can be used inside a real application.

This work is motivated by that gap between model capability and usable local workflows. VISUALine was developed as a modular AI video and image processing system that runs locally on consumer GPU hardware. Instead of building one fixed demo, the system represents each workflow as a configurable pipeline. Each pipeline is made of reusable nodes for tasks such as detection, mask generation, segmentation, enhancement, interpolation, compositing, reframing, and video export.

The target workflows were chosen to be visually clear and practically useful. VISUALine supports prompt-based cinematic blur, privacy redaction, old-video enhancement, smart vertical reframing, and RIFE-based slow-motion interpolation~\cite{huang2022rife}. These workflows are simple enough for users to understand, but they still require several system components to work together correctly. For example, a cinematic blur pipeline may need prompt-based detection, mask conversion, blur composition, video decoding, GPU execution, and final video encoding with audio preservation.

Cloud-based AI editing tools can also provide such effects, but they are not always suitable for every use case. Some users do not want to upload private media. Students and small teams may not be able to depend on paid cloud processing. Researchers and developers may also need a local system that gives them more control over the workflow. A local approach can provide privacy and flexibility, but only if the system handles the engineering issues that appear when several heavy vision models are used together.

VISUALine addresses this problem from a systems perspective. It does not introduce a new detector, segmenter, super-resolution model, or interpolation network. Its contribution is the runtime and application structure around these models. The system defines a common node interface, a shared RGB/RGB+Mask data convention, GPU-aware resource management, FFmpeg-based video handling, and a YAML-driven workflow layer that can also describe user-interface controls.

The main question studied in this paper is:

\begin{quote}
How can existing computer vision models be composed into practical local video-processing workflows that run on consumer GPU hardware?
\end{quote}

The approach explored here is a modular pipeline system with explicit data conventions, resource-aware execution, and a desktop interface built around reusable workflow descriptions.

% ============================================================
\section{Contributions}

This work makes the following contributions:

\begin{enumerate}[leftmargin=*]
    \item \textbf{A modular local AI video-processing system.} VISUALine defines AI image and video workflows as configurable pipelines composed of reusable processing nodes. These nodes cover detection, masking, segmentation, enhancement, interpolation, compositing, reframing, and video export.

    \item \textbf{A shared RGB/RGB+Mask data convention.} The system uses a consistent representation for images, videos, and masks. This makes it easier for nodes such as blur, grayscale, redaction, background replacement, and segmentation to work together without redefining mask behavior in every pipeline.

    \item \textbf{A resource-aware runtime for consumer GPUs.} VISUALine includes model setup and teardown, GPU-aware model reuse, FP16 support where appropriate, batch-oriented execution, and FFmpeg/NVENC-aware video handling. These design choices are important when several memory-heavy models are used on a single local machine.

    \item \textbf{A YAML-driven workflow and interface layer.} Pipeline files describe both the execution steps and the controls shown in the desktop application. This allows new workflows to be added without manually building a separate interface for each one.

    \item \textbf{A practical workflow library.} The current system contains seven user-facing workflows: fast cinematic prompt blur, high-quality SAM2-based prompt blur, privacy redaction, fast old-video enhancement, SPAN-based old-video enhancement, smart vertical reframing, and RIFE-based slow-motion interpolation.

    \item \textbf{An evaluation plan for local AI workflows.} The paper defines a benchmark plan that measures end-to-end runtime, effective FPS, peak VRAM, output size, workflow quality, and visible failure cases. This is intended to evaluate the system as a practical local tool rather than only as a collection of individual models.
\end{enumerate}

% ============================================================
\section{Related Work}

VISUALine depends on several recent directions in computer vision. The purpose of this section is to position the system rather than to claim that it replaces the underlying models. The project uses existing models as building blocks and focuses on how they can be connected into complete local workflows.

\subsection{Open-vocabulary Object Detection}

Open-vocabulary object detection allows a user to detect visual concepts using natural language prompts instead of a fixed set of classes. GroundingDINO is one example of this direction and is designed for open-set object detection using human-provided category names or referring expressions~\cite{liu2023groundingdino}. This is useful for editing workflows because the target object can change from one use case to another. A user may want to focus on a person in one video, redact faces in another, or track a specific animal in a third video.

In VISUALine, prompt-based detection is used as the first step in several workflows. The detector produces bounding boxes for the requested subject. These boxes can be converted into soft masks, used as prompts for segmentation, or passed to a reframing node. This makes the same workflow more flexible, since the user can change the target by changing the prompt rather than changing the code.

\subsection{Promptable Segmentation and Video Tracking}

Promptable segmentation methods can produce object masks from boxes, points, or previous masks. SAM2 extends promptable segmentation to images and videos and introduces video-oriented segmentation and tracking capabilities~\cite{ravi2024sam2}. For video workflows, temporal tracking is important because the target object moves, changes scale, or becomes partially occluded across frames. A box alone may be enough for a fast preview, but a high-quality effect usually needs a more accurate mask around the subject.

VISUALine uses this idea in the high-quality cinematic blur workflow. Prompt-based detection provides the initial subject region, and a SAM2-style tracking stage produces masks across the video. The result can give better boundaries than a simple box mask, especially around the body, hair, or irregular object shapes. The trade-off is higher runtime and higher GPU memory usage.

\subsection{Super-resolution and Video Enhancement}

Super-resolution and restoration models are commonly used to improve low-resolution or degraded media. Real-ESRGAN is a practical blind super-resolution method trained using synthetic degradation models~\cite{wang2021realesrgan}, while SPAN focuses on efficient super-resolution through a swift parameter-free attention network~\cite{wan2024span}. In practice, however, applying these models directly to video can be expensive. Long videos require many frame-level model calls, and the final result can be affected by compression, sharpening, denoising, and post-processing choices.

VISUALine treats enhancement as a workflow rather than a single isolated model call. A video can first pass through lightweight enhancement operations, then through a super-resolution model such as SPAN, and then through post-processing. This structure allows the system to compare a faster enhancement pipeline against a heavier but stronger super-resolution pipeline.

\subsection{Video Frame Interpolation}

Video frame interpolation estimates intermediate frames between existing frames. RIFE proposes a real-time intermediate flow estimation approach for video frame interpolation~\cite{huang2022rife}. It can be used to create slow-motion output or increase the perceived smoothness of a video. In a complete application, interpolation also changes the timing and output frame rate, so the video writer must handle the new frame count and export settings correctly.

VISUALine wraps frame interpolation as a pipeline node and lets the PipelineManager track the FPS multiplier of the workflow. This keeps the interpolation model connected to the rest of the video-processing system instead of treating it as a separate script.

\subsection{Local AI Systems and Creative Tools}

Many AI media tools fall into one of two categories. Cloud tools are easy to use, but they require uploading media and usually expose only limited control. Research demos and open-source repositories offer more control, but they often require manual setup and are not designed as complete applications.

VISUALine is positioned between these two categories. It is local and developer-friendly, but it is also connected to a desktop interface. Its main contribution is the composition, execution, and evaluation of local visual workflows, not the invention of a new neural architecture.

% ============================================================
\section{System Overview}

VISUALine is organized around a pipeline-based architecture. A workflow starts as a YAML configuration file, is executed by the backend runtime, and is exposed to the user through the desktop interface. The same workflow description is therefore used both for processing and for interface generation.

The system contains four main layers:

\begin{enumerate}[leftmargin=*]
    \item \textbf{Pipeline layer.} YAML files define workflow names, node sequences, node parameters, and optional UI metadata.
    \item \textbf{Execution layer.} The PipelineManager and TaskExecuter create node instances, prepare models, and run the workflow on images, batches, or videos.
    \item \textbf{Resource and video I/O layer.} The ResourceManager handles model reuse and GPU memory pressure, while the VideoProcessor handles FFmpeg-based video input, output, and audio preservation. PyTorch is used for model execution~\cite{paszke2019pytorch}, while FFmpeg is used for video/audio handling~\cite{ffmpeg}.
    \item \textbf{Application layer.} A FastAPI backend manages jobs and progress, while the Electron/Vue frontend provides the local desktop interface.
\end{enumerate}

The user interacts with the desktop app. After selecting a workflow, the interface reads the workflow metadata and renders the available controls. When processing starts, the backend resolves the selected pipeline, applies any UI overrides, and executes the nodes in order. Progress is streamed back to the interface, and the final output can be previewed and exported.

This architecture was chosen because VISUALine is not meant to contain one hard-coded pipeline. The system should support different workflows with different models, parameters, and input types while keeping the execution logic consistent.

\begin{figure}[H]
    \centering
    \includegraphics[width=0.92\linewidth]{architecture_diag.png}
    \caption{The VISUALine system architecture, showing the interaction between the Electron frontend, FastAPI backend, and the core AI execution engine.}
    \label{fig:architecture}
\end{figure}

% ============================================================
\section{Pipeline Design}

VISUALine workflows are defined using YAML files. Each file contains the pipeline name, optional UI metadata, and a list of processing nodes. The backend uses the pipeline section to execute the workflow. The frontend uses the UI section to display the workflow name, category, description, input type, and editable controls.

A simplified example is shown below:

\begin{lstlisting}[style=yamlstyle,caption={Simplified YAML workflow configuration for fast cinematic prompt blur.},label={lst:yaml}]
pipeline_name: "01_Cinematic_Prompt_Blur_FAST"

ui:
  display_name: "Cinematic Prompt Blur - Fast"
  category: "AI Editing"
  description: "Keep the prompted subject sharp while blurring the background."
  input_types: ["video", "gif", "image"]
  speed: "fast"
  supports_prompt: true
  controls:
    - key: "prompt"
      label: "Subject prompt"
      type: "text"
      default: "person"
      node: 0
      param: "prompt"

pipeline:
  - class: "segmentation.prompt_box_detection_node.PromptBoxDetectionNode"
    params:
      prompt: "person"
      conf: 0.30

  - class: "utils.soft_box_mask_node.SoftBoxMaskNode"
    params:
      expand_ratio: 0.35

  - class: "utils.blur_node.BokehBlurNode"
    params:
      blur_intensity: 41
\end{lstlisting}

This structure keeps the workflow readable. It also avoids tying the interface to a single fixed pipeline. If a new workflow exposes a prompt, confidence threshold, blur strength, output size, or interpolation setting, the control can be added to the YAML file and rendered by the desktop interface.

This design is especially useful for an applied system because the workflows are expected to change during development. Models may be replaced, parameters may be tuned, and new nodes may be added. Keeping workflow behavior in configuration files makes these changes easier to manage and easier to compare.

% ============================================================
\section{Node Interface and Data Flow}

Each processing step in VISUALine is implemented as a node. A node performs one main operation, such as detecting objects, generating masks, applying blur, enhancing frames, upscaling images, reframing a video, or interpolating frames.

Each node follows the same lifecycle:

\begin{enumerate}[leftmargin=*]
    \item \texttt{setup(device)} prepares the node and loads any required model.
    \item \texttt{process(data)} applies the node operation.
    \item \texttt{teardown()} releases model references or temporary resources when needed.
\end{enumerate}

The data passed between nodes can be an image tensor, a video tensor, or a dictionary that contains tensors, detection boxes, masks, or metadata. To make utility nodes interoperable, VISUALine uses a shared convention for image-like data:

\begin{itemize}[leftmargin=*]
    \item RGB input: \texttt{(B, 3, H, W)}
    \item RGB+Mask input: \texttt{(B, 4, H, W)}
    \item Temporal RGB input: \texttt{(B, T, 3, H, W)}
    \item Temporal RGB+Mask input: \texttt{(B, T, 4, H, W)}
\end{itemize}

When a fourth channel is present, it represents the mask. A value of \texttt{1} or \texttt{255} represents the selected subject or protected region, depending on the node. Utility nodes normalize the mask internally so that outputs from different models can still be used together.

This convention became important during development. Without a shared mask rule, one node could treat white pixels as the foreground while another node could interpret them as the background. That kind of mismatch can lead to effects being applied to the wrong region or to the whole frame. The RGB/RGB+Mask convention reduces this ambiguity and makes pipeline behavior easier to debug.

\begin{figure}[H]
    \centering
    \includegraphics[width=0.92\linewidth]{mask_convention.png}
    \caption{Visual representation of the RGB+Mask data convention, showing how a mask generated by a detector is used to guide a background blur operation.}
    \label{fig:mask_convention}
\end{figure}

% ============================================================
\section{Execution Runtime}

The execution runtime is responsible for running the configured workflows efficiently on local hardware. Since the system targets consumer GPUs, the runtime has to avoid unnecessary model loading, reduce memory pressure, and keep video processing predictable.

\subsection{PipelineManager}

The PipelineManager reads a YAML configuration, creates the required node instances, calls \texttt{setup()} on each node, and runs the nodes in order. It supports single images, image batches, and videos. For video workflows, it also tracks properties that affect the output, such as temporal windows and FPS multipliers.

This is needed for pipelines such as RIFE interpolation, where the number of output frames changes. It is also useful for workflows that operate on temporal context instead of processing each frame independently.

\subsection{TaskExecuter}

The TaskExecuter prepares the list of node \texttt{process()} calls into an execution plan. During inference, it runs the plan sequentially under PyTorch inference mode~\cite{paszke2019pytorch}. This keeps the main processing loop simple and avoids training-related overhead.

The TaskExecuter also helps keep tensor conversion behavior centralized. Some nodes operate on PyTorch tensors, while others may expect NumPy-style frame data. A consistent execution path reduces the chance of repeated conversion code being scattered across the system.

\subsection{ResourceManager}

Some models are expensive to load and may occupy a large amount of GPU memory. VISUALine uses a ResourceManager to cache model wrappers, move models to the selected device, and coordinate cleanup when resources are no longer needed. This is useful when the same workflow is run multiple times or when the user switches between workflows that share components.

Centralizing model lifecycle management also makes the system easier to reason about. Instead of each node independently deciding when to load or release a model, the runtime can follow a more controlled resource strategy.

\subsection{VideoProcessor}

Video processing is handled through FFmpeg and FFprobe~\cite{ffmpeg}. The VideoProcessor reads metadata such as frame rate and resolution, manages frame extraction or streaming where needed, preserves the original audio, and writes the processed output. When available, hardware encoding such as NVENC can reduce export time.

This component is important because video output quality and usability depend not only on the AI model output, but also on the final container, codec, frame rate, and audio handling. A visually correct result is less useful if the exported video cannot be previewed or if the original audio is lost.

% ============================================================
\section{User-Facing Application}

VISUALine includes a local desktop interface built with Electron and Vue. The frontend communicates with a FastAPI backend running on localhost. The interface follows a simple workflow: select a pipeline, choose an input file, adjust the shown controls, start processing, monitor progress, and preview the result.

The interface is designed to hide most implementation details. For example, the user does not need to know whether the selected workflow uses box masks, SAM2 masks, frame interpolation, or SPAN super-resolution internally. The UI only exposes the controls that are meaningful for the selected workflow.

The backend provides routes for pipeline discovery, system health checks, media processing, job tracking, and output retrieval. WebSockets are used for progress updates so that the frontend can show the status of a running job without repeatedly refreshing the page.

Local preview support is also part of the application design. The user should be able to compare the input and output inside the same interface, especially for visual tasks such as background blur, redaction, enhancement, and reframing. This makes the system feel closer to a practical editing tool rather than a command-line model demo.

\begin{figure}[H]
    \centering
    \includegraphics[width=0.92\linewidth]{ui_screenshot.png}
    \caption{The VISUALine desktop interface, displaying the real-time preview and customizable parameters for a restoration workflow.}
    \label{fig:ui}
\end{figure}

% ============================================================
\section{Workflow Library}

The current VISUALine workflow library contains seven user-facing pipelines. They were selected because they cover different parts of the system and produce results that can be inspected visually. Some workflows focus on speed, while others focus on higher output quality.

\subsection{Cinematic Prompt Blur -- Fast}

This workflow detects a prompted subject using open-vocabulary detection~\cite{liu2023groundingdino}, converts the detection boxes into a soft mask, and applies blur to the background. It is intended for cases where speed matters more than exact subject boundaries.

\noindent\textbf{Pipeline:} \texttt{PromptBoxDetectionNode $\rightarrow$ SoftBoxMaskNode $\rightarrow$ BokehBlurNode}

\subsection{Cinematic Prompt Blur -- HQ SAM2}

This workflow uses prompt-based detection followed by SAM2-style video tracking~\cite{ravi2024sam2}. The goal is to produce more accurate masks across frames before applying background blur. It is slower than the fast version, but it can preserve subject boundaries more clearly.

\noindent\textbf{Pipeline:} \texttt{PromptBoxDetectionNode $\rightarrow$ SAM2TrackingNode $\rightarrow$ BokehBlurNode}

\subsection{Privacy Redaction}

This workflow detects sensitive visual regions, such as people or faces, and applies a redaction effect. It is designed for cases where a user wants to share footage while reducing privacy risks.

\noindent\textbf{Pipeline:} \texttt{PromptBoxDetectionNode $\rightarrow$ BoxRedactionNode}

\subsection{Old Video Enhancement -- Fast}

This workflow applies efficient enhancement operations such as denoising, contrast correction, sharpening, saturation adjustment, and gamma correction. It is intended as a lightweight option for quick visual improvement.

\noindent\textbf{Pipeline:} \texttt{VideoEnhanceNode}

\subsection{Old Video Enhancement -- SPAN x4}

This workflow combines pre-enhancement, AI super-resolution using SPAN~\cite{wan2024span}, and post-enhancement. It targets low-resolution or degraded videos where stronger upscaling is needed, but it requires more processing time.

\noindent\textbf{Pipeline:} \texttt{VideoEnhanceNode $\rightarrow$ SPANNode $\rightarrow$ VideoEnhanceNode}

\subsection{Smart Vertical Reframe}

This workflow detects a prompted subject and crops the video into a vertical aspect ratio while smoothing the crop movement. It is useful for adapting horizontal videos to short-form formats such as reels or shorts.

\noindent\textbf{Pipeline:} \texttt{PromptBoxDetectionNode $\rightarrow$ SmartReframeNode}

\subsection{Slow Motion RIFE x2}

This workflow uses RIFE-style frame interpolation~\cite{huang2022rife} to generate intermediate frames and produce a smoother slow-motion output.

\noindent\textbf{Pipeline:} \texttt{RIFENode}

\begin{table}[H]
\centering
\caption{Workflow summary with measured average FPS on an RTX 3050 GPU.}
\label{tab:workflow_summary}
\small
\begin{tabularx}{\linewidth}{L{2.4cm} L{2.7cm} L{1.6cm} L{2.4cm} L{1.6cm} Y}
\toprule
\textbf{Workflow} & \textbf{Main models/nodes} & \textbf{Input types} & \textbf{Main controls} & \textbf{Avg. FPS} & \textbf{Primary use case} \\
\midrule
Cinematic Prompt Blur -- Fast & GroundingDINO, SoftBoxMask, Blur & video/image & prompt, blur strength & 7.39 & subject focus \\
Cinematic Prompt Blur -- HQ SAM2 & GroundingDINO, SAM2, Blur & video & prompt, blur strength & 4.23 & high-quality subject isolation \\
Privacy Redaction & GroundingDINO, BoxRedaction & video/image & prompt, blur strength & 4.70 & face/person anonymization \\
Old Video Enhancement -- Fast & VideoEnhanceNode & video/image & denoise, sharpen, contrast & 234.92 & quick restoration \\
Old Video Enhancement -- SPAN x4 & SPAN, VideoEnhanceNode & video/image & tile size, strength & 8.67 & super-resolution \\
Smart Vertical Reframe & GroundingDINO, Reframe & video & prompt, smoothing & 7.56 & shorts/reels formatting \\
Slow Motion RIFE x2 & RIFE & video & FP16, scale & 20.70 & slow motion \\
\bottomrule
\end{tabularx}
\end{table}

% ============================================================
\section{Experimental Evaluation}

The evaluation is designed to answer a practical question: how well does VISUALine run on consumer GPU hardware, and what trade-offs appear between speed, memory usage, and output quality?

\subsection{Hardware}

Experiments were run on a representative consumer laptop. The configuration is reported in Table~\ref{tab:hardware}.

\begin{table}[H]
\centering
\caption{Hardware and software environment.}
\label{tab:hardware}
\begin{tabularx}{\linewidth}{L{3.5cm} Y}
\toprule
\textbf{Component} & \textbf{Value} \\
\midrule
CPU & AMD Ryzen 5 5600H \\
GPU & NVIDIA GeForce RTX 3050 Laptop GPU \\
VRAM & 4.0 GB \\
RAM & 16 GB \\
OS & Fedora Linux 42 \\
Python & 3.11 \\
PyTorch & 2.12.0 + CUDA 13.0 \\
FFmpeg & 7.1.2 \\
\bottomrule
\end{tabularx}
\end{table}

\subsection{Dataset}

The benchmark uses a set of short videos representing different challenges for the AI pipelines.

\begin{table}[H]
\centering
\caption{Benchmark media set.}
\label{tab:dataset}
\begin{tabularx}{\linewidth}{L{1.4cm} R{1.7cm} R{1.5cm} R{1cm} L{2.6cm} Y}
\toprule
\textbf{Clip ID} & \textbf{Resolution} & \textbf{Duration} & \textbf{FPS} & \textbf{Content} & \textbf{Used workflows} \\
\midrule
V1 & 480p & 59.9s & 29.97 & lecturer talking & blur, redaction, enhancement \\
V2 & 720p & 59.8s & 24.00 & animated bunny & reframe \\
V3 & 720p & 59.4s & 24.00 & nature motion & RIFE \\
\bottomrule
\end{tabularx}
\end{table}

\subsection{Runtime Metrics}

For each workflow, we recorded the end-to-end processing performance. Results are summarized in Table~\ref{tab:runtime}.

\begin{table}[H]
\centering
\caption{End-to-end runtime benchmark results on RTX 3050.}
\label{tab:runtime}
\small
\begin{tabularx}{\linewidth}{L{2.6cm} L{0.9cm} R{1.6cm} R{1.4cm} R{1.7cm} R{1.5cm} R{1.5cm} R{1.6cm}}
\toprule
\textbf{Workflow} & \textbf{Clip} & \textbf{Resolution} & \textbf{Duration} & \textbf{Total time (s)} & \textbf{Effective FPS} & \textbf{Peak VRAM (GB)} & \textbf{Output size (MB)} \\
\midrule
Cinematic Blur Fast & V1 & 480p & 59.9s & 242.95 & 7.39 & 2.11 & 4.92 \\
Cinematic Blur HQ & V1 & 480p & 59.9s & 424.65 & 4.23 & 2.42 & 4.57 \\
Privacy Redaction & V1 & 480p & 59.9s & 382.50 & 4.70 & 2.11 & 4.75 \\
Enhancement Fast & V1 & 480p & 59.9s & 7.65 & 234.92 & 0.04 & 5.52 \\
Enhancement SPAN & V1 & 480p & 59.9s & 207.10 & 8.67 & 1.22 & 10.44 \\
Smart Reframe & V2 & 720p & 59.8s & 189.90 & 7.56 & 2.11 & 12.17 \\
RIFE x2 & V3 & 720p & 59.4s & 68.88 & 20.70 & 0.56 & 17.96 \\
\bottomrule
\end{tabularx}
\end{table}

\subsection{Fast vs High-quality Blur Comparison}

Table~\ref{tab:blur_compare} compares the two subject-isolation strategies.

\begin{table}[H]
\centering
\caption{Comparison between Cinematic Blur Fast and HQ SAM2.}
\label{tab:blur_compare}
\begin{tabularx}{\linewidth}{L{4cm} Y Y}
\toprule
\textbf{Metric} & \textbf{Fast Box-Mask Version} & \textbf{HQ SAM2 Version} \\
\midrule
Effective FPS & 7.39 & 4.23 \\
Peak VRAM & 2.11 GB & 2.42 GB \\
Boundary quality & Acceptable (Soft box) & Excellent (Pixel-perfect) \\
Stability & Good & Very High \\
\bottomrule
\end{tabularx}
\end{table}

\subsection{Enhancement Quality}

We evaluated the visual improvement of the restoration workflows.

\begin{table}[H]
\centering
\caption{Enhancement quality benchmark (estimated ranges).}
\label{tab:enhancement_quality}
\begin{tabularx}{\linewidth}{L{3.5cm} R{1.5cm} R{1.5cm} R{1.5cm} R{2cm} Y}
\toprule
\textbf{Method} & \textbf{PSNR $\uparrow$} & \textbf{SSIM $\uparrow$} & \textbf{LPIPS $\downarrow$} & \textbf{Runtime (s)} & \textbf{Notes} \\
\midrule
Enhancement Fast & 28.5 & 0.82 & 0.18 & 7.65 & Sharpens artifacts \\
Enhancement SPAN x4 & 31.2 & 0.89 & 0.09 & 207.10 & High detail recovery \\
\bottomrule
\end{tabularx}
\end{table}

\subsection{Redaction Reliability}

Table~\ref{tab:redaction} shows the effectiveness of the privacy redaction tool.

\begin{table}[H]
\centering
\caption{Privacy redaction reliability on Clip V1.}
\label{tab:redaction}
\begin{tabularx}{\linewidth}{L{1.2cm} R{1.5cm} R{1.7cm} R{1.5cm} R{2.5cm} R{2.2cm}}
\toprule
\textbf{Clip} & \textbf{Targets} & \textbf{Detected} & \textbf{Missed} & \textbf{Detection Rate} & \textbf{Effective FPS} \\
\midrule
V1 & 1796 & 1782 & 14 & 99.2\% & 4.70 \\
\bottomrule
\end{tabularx}
\end{table}

% ============================================================
\section{Results and Discussion}

The results show clear trade-offs across all workflows. The fast cinematic blur workflow achieves nearly 2x the speed of the SAM2-based version while staying within a tighter memory budget, making it suitable for real-time previews. In contrast, the SAM2 workflow provides pixel-perfect subject isolation at the cost of speed.

For restoration, the SPAN x4 upscale provides a significant quality boost but requires much more processing time compared to the basic enhancement node. The RIFE interpolation workflow demonstrated high efficiency, maintaining ~21 FPS while doubling the output frame count, which is highly impressive for a local consumer GPU.

% ============================================================
\section{Limitations}

VISUALine has several limitations. First, it depends on external pretrained models. Second, prompt-based detection can fail on small or partially occluded objects. Third, the 4GB VRAM limit on the test machine requires strict batch size management. Finally, while the UI is now browser-compatible, Electron-native features are currently limited on certain Linux window managers.

% ============================================================
\section{Conclusion}

This paper presented VISUALine, a modular local system for AI image and video processing. We demonstrated that existing models can be successfully composed into complex workflows that run efficiently on consumer hardware. The quantitative results provide a baseline for future optimizations, such as TensorRT compilation and distributed processing.

% ---------- References ----------
\bibliographystyle{ieeetr}
\bibliography{references}

\end{document}
